% Options for packages loaded elsewhere
\PassOptionsToPackage{unicode}{hyperref}
\PassOptionsToPackage{hyphens}{url}
%
\documentclass[
  9pt,
  ignorenonframetext,
]{beamer}
\usepackage{pgfpages}
\setbeamertemplate{caption}[numbered]
\setbeamertemplate{caption label separator}{: }
\setbeamercolor{caption name}{fg=normal text.fg}
\beamertemplatenavigationsymbolsempty
% Prevent slide breaks in the middle of a paragraph
\widowpenalties 1 10000
\raggedbottom
\setbeamertemplate{part page}{
  \centering
  \begin{beamercolorbox}[sep=16pt,center]{part title}
    \usebeamerfont{part title}\insertpart\par
  \end{beamercolorbox}
}
\setbeamertemplate{section page}{
  \centering
  \begin{beamercolorbox}[sep=12pt,center]{part title}
    \usebeamerfont{section title}\insertsection\par
  \end{beamercolorbox}
}
\setbeamertemplate{subsection page}{
  \centering
  \begin{beamercolorbox}[sep=8pt,center]{part title}
    \usebeamerfont{subsection title}\insertsubsection\par
  \end{beamercolorbox}
}
\AtBeginPart{
  \frame{\partpage}
}
\AtBeginSection{
  \ifbibliography
  \else
    \frame{\sectionpage}
  \fi
}
\AtBeginSubsection{
  \frame{\subsectionpage}
}

\usepackage{amsmath,amssymb}
\usepackage{iftex}
\ifPDFTeX
  \usepackage[T1]{fontenc}
  \usepackage[utf8]{inputenc}
  \usepackage{textcomp} % provide euro and other symbols
\else % if luatex or xetex
  \usepackage{unicode-math}
  \defaultfontfeatures{Scale=MatchLowercase}
  \defaultfontfeatures[\rmfamily]{Ligatures=TeX,Scale=1}
\fi
\usepackage{lmodern}
\usetheme[]{Copenhagen}
\ifPDFTeX\else  
    % xetex/luatex font selection
\fi
% Use upquote if available, for straight quotes in verbatim environments
\IfFileExists{upquote.sty}{\usepackage{upquote}}{}
\IfFileExists{microtype.sty}{% use microtype if available
  \usepackage[]{microtype}
  \UseMicrotypeSet[protrusion]{basicmath} % disable protrusion for tt fonts
}{}
\makeatletter
\@ifundefined{KOMAClassName}{% if non-KOMA class
  \IfFileExists{parskip.sty}{%
    \usepackage{parskip}
  }{% else
    \setlength{\parindent}{0pt}
    \setlength{\parskip}{6pt plus 2pt minus 1pt}}
}{% if KOMA class
  \KOMAoptions{parskip=half}}
\makeatother
\usepackage{xcolor}
\newif\ifbibliography
\setlength{\emergencystretch}{3em} % prevent overfull lines
\setcounter{secnumdepth}{-\maxdimen} % remove section numbering

\usepackage{color}
\usepackage{fancyvrb}
\newcommand{\VerbBar}{|}
\newcommand{\VERB}{\Verb[commandchars=\\\{\}]}
\DefineVerbatimEnvironment{Highlighting}{Verbatim}{commandchars=\\\{\}}
% Add ',fontsize=\small' for more characters per line
\usepackage{framed}
\definecolor{shadecolor}{RGB}{241,243,245}
\newenvironment{Shaded}{\begin{snugshade}}{\end{snugshade}}
\newcommand{\AlertTok}[1]{\textcolor[rgb]{0.68,0.00,0.00}{#1}}
\newcommand{\AnnotationTok}[1]{\textcolor[rgb]{0.37,0.37,0.37}{#1}}
\newcommand{\AttributeTok}[1]{\textcolor[rgb]{0.40,0.45,0.13}{#1}}
\newcommand{\BaseNTok}[1]{\textcolor[rgb]{0.68,0.00,0.00}{#1}}
\newcommand{\BuiltInTok}[1]{\textcolor[rgb]{0.00,0.23,0.31}{#1}}
\newcommand{\CharTok}[1]{\textcolor[rgb]{0.13,0.47,0.30}{#1}}
\newcommand{\CommentTok}[1]{\textcolor[rgb]{0.37,0.37,0.37}{#1}}
\newcommand{\CommentVarTok}[1]{\textcolor[rgb]{0.37,0.37,0.37}{\textit{#1}}}
\newcommand{\ConstantTok}[1]{\textcolor[rgb]{0.56,0.35,0.01}{#1}}
\newcommand{\ControlFlowTok}[1]{\textcolor[rgb]{0.00,0.23,0.31}{\textbf{#1}}}
\newcommand{\DataTypeTok}[1]{\textcolor[rgb]{0.68,0.00,0.00}{#1}}
\newcommand{\DecValTok}[1]{\textcolor[rgb]{0.68,0.00,0.00}{#1}}
\newcommand{\DocumentationTok}[1]{\textcolor[rgb]{0.37,0.37,0.37}{\textit{#1}}}
\newcommand{\ErrorTok}[1]{\textcolor[rgb]{0.68,0.00,0.00}{#1}}
\newcommand{\ExtensionTok}[1]{\textcolor[rgb]{0.00,0.23,0.31}{#1}}
\newcommand{\FloatTok}[1]{\textcolor[rgb]{0.68,0.00,0.00}{#1}}
\newcommand{\FunctionTok}[1]{\textcolor[rgb]{0.28,0.35,0.67}{#1}}
\newcommand{\ImportTok}[1]{\textcolor[rgb]{0.00,0.46,0.62}{#1}}
\newcommand{\InformationTok}[1]{\textcolor[rgb]{0.37,0.37,0.37}{#1}}
\newcommand{\KeywordTok}[1]{\textcolor[rgb]{0.00,0.23,0.31}{\textbf{#1}}}
\newcommand{\NormalTok}[1]{\textcolor[rgb]{0.00,0.23,0.31}{#1}}
\newcommand{\OperatorTok}[1]{\textcolor[rgb]{0.37,0.37,0.37}{#1}}
\newcommand{\OtherTok}[1]{\textcolor[rgb]{0.00,0.23,0.31}{#1}}
\newcommand{\PreprocessorTok}[1]{\textcolor[rgb]{0.68,0.00,0.00}{#1}}
\newcommand{\RegionMarkerTok}[1]{\textcolor[rgb]{0.00,0.23,0.31}{#1}}
\newcommand{\SpecialCharTok}[1]{\textcolor[rgb]{0.37,0.37,0.37}{#1}}
\newcommand{\SpecialStringTok}[1]{\textcolor[rgb]{0.13,0.47,0.30}{#1}}
\newcommand{\StringTok}[1]{\textcolor[rgb]{0.13,0.47,0.30}{#1}}
\newcommand{\VariableTok}[1]{\textcolor[rgb]{0.07,0.07,0.07}{#1}}
\newcommand{\VerbatimStringTok}[1]{\textcolor[rgb]{0.13,0.47,0.30}{#1}}
\newcommand{\WarningTok}[1]{\textcolor[rgb]{0.37,0.37,0.37}{\textit{#1}}}

\providecommand{\tightlist}{%
  \setlength{\itemsep}{0pt}\setlength{\parskip}{0pt}}\usepackage{longtable,booktabs,array}
\usepackage{calc} % for calculating minipage widths
\usepackage{caption}
% Make caption package work with longtable
\makeatletter
\def\fnum@table{\tablename~\thetable}
\makeatother
\usepackage{graphicx}
\makeatletter
\def\maxwidth{\ifdim\Gin@nat@width>\linewidth\linewidth\else\Gin@nat@width\fi}
\def\maxheight{\ifdim\Gin@nat@height>\textheight\textheight\else\Gin@nat@height\fi}
\makeatother
% Scale images if necessary, so that they will not overflow the page
% margins by default, and it is still possible to overwrite the defaults
% using explicit options in \includegraphics[width, height, ...]{}
\setkeys{Gin}{width=\maxwidth,height=\maxheight,keepaspectratio}
% Set default figure placement to htbp
\makeatletter
\def\fps@figure{htbp}
\makeatother

\makeatletter
\@ifpackageloaded{caption}{}{\usepackage{caption}}
\AtBeginDocument{%
\ifdefined\contentsname
  \renewcommand*\contentsname{Table of contents}
\else
  \newcommand\contentsname{Table of contents}
\fi
\ifdefined\listfigurename
  \renewcommand*\listfigurename{List of Figures}
\else
  \newcommand\listfigurename{List of Figures}
\fi
\ifdefined\listtablename
  \renewcommand*\listtablename{List of Tables}
\else
  \newcommand\listtablename{List of Tables}
\fi
\ifdefined\figurename
  \renewcommand*\figurename{Figure}
\else
  \newcommand\figurename{Figure}
\fi
\ifdefined\tablename
  \renewcommand*\tablename{Table}
\else
  \newcommand\tablename{Table}
\fi
}
\@ifpackageloaded{float}{}{\usepackage{float}}
\floatstyle{ruled}
\@ifundefined{c@chapter}{\newfloat{codelisting}{h}{lop}}{\newfloat{codelisting}{h}{lop}[chapter]}
\floatname{codelisting}{Listing}
\newcommand*\listoflistings{\listof{codelisting}{List of Listings}}
\makeatother
\makeatletter
\makeatother
\makeatletter
\@ifpackageloaded{caption}{}{\usepackage{caption}}
\@ifpackageloaded{subcaption}{}{\usepackage{subcaption}}
\makeatother

\ifLuaTeX
  \usepackage{selnolig}  % disable illegal ligatures
\fi
\usepackage{bookmark}

\IfFileExists{xurl.sty}{\usepackage{xurl}}{} % add URL line breaks if available
\urlstyle{same} % disable monospaced font for URLs
\hypersetup{
  pdftitle={R básico para ciencia de datos - Clase 1},
  hidelinks,
  pdfcreator={LaTeX via pandoc}}


\title{R básico para ciencia de datos - Clase 1}
\author{}
\date{Invalid Date}

\begin{document}
\frame{\titlepage}


\begin{frame}{Generalidades de R --- Día 1}
\phantomsection\label{generalidades-de-r-duxeda-1}
\begin{itemize}
\tightlist
\item
  \textbf{¿Qué es R y por qué usarlo?}
\item
  \textbf{RStudio y paneles}
\item
  \textbf{Configuración del ambiente de trabajo y Scripts}
\item
  \textbf{Directorios y sesión}
\item
  \textbf{Áreas de trabajo y Proyectos}
\item
  \textbf{Instalación y carga de paquetes}
\item
  \textbf{Tipos de objetos, tipos de datos y estructuras de datos}
\item
  \textbf{Funciones}
\end{itemize}
\end{frame}

\begin{frame}{¿Qué es R?}
\phantomsection\label{quuxe9-es-r}
\begin{itemize}
\tightlist
\item
  R es un lenguaje de programación \textbf{enfocado en estadística y
  análisis de datos}
\item
  Fue creado por estadísticos como un ambiente interactivo para el
  análisis de datos
\item
  Puede guardar su trabajo como una secuencia de comandos
  (\emph{script}) que se puede ejecutar fácilmente en cualquier momento
\item
  Otros programas (C, Java, Python) también pueden hacer cálculos
  estadísticos, pero R fue diseñado específicamente para eso
\end{itemize}
\end{frame}

\begin{frame}{Un poco de historia\ldots{}}
\phantomsection\label{un-poco-de-historia}
\begin{itemize}
\tightlist
\item
  R proviene del lenguaje \textbf{S}, creado en los Laboratorios Bell
  (EE.UU.) --- los mismos que inventaron el transistor, el láser y Unix
\item
  \textbf{Ross Ihaka} y \textbf{Robert Gentleman} (Universidad de
  Auckland, NZ) crearon una implementación libre y gratuita de S
\item
  El trabajo inició en \textbf{1992} y la primera versión estable llegó
  en \textbf{2000}
\item
  Hoy es mantenido por el \emph{R Development Core Team}, con
  especialistas de todo el mundo
\item
  R es software libre con \textbf{Licencia Pública GNU} --- puedes
  usarlo para fines personales, académicos o comerciales sin
  restricciones
\end{itemize}
\end{frame}

\begin{frame}{¿Por qué usar R?}
\phantomsection\label{por-quuxe9-usar-r}
\begin{enumerate}
\tightlist
\item
  Gratuito y de código abierto
\item
  Corre en Windows, Mac OS y UNIX/Linux --- scripts compatibles entre
  plataformas
\item
  Comunidad grande, creciente y activa
\item
  Acceso temprano a nuevas metodologías publicadas como paquetes
\item
  Miles de paquetes disponibles en \textbf{CRAN}, \textbf{GitHub} y
  \textbf{Bioconductor} para: ecología, biología molecular, ciencias
  sociales, geografía, estadística\ldots{}
\end{enumerate}
\end{frame}

\begin{frame}{RStudio}
\phantomsection\label{rstudio}
RStudio es nuestra plataforma de trabajo. Provee un editor visual e
interactivo para crear y editar \emph{scripts}, además de otras
herramientas útiles.

\begin{center}
\includegraphics[width=3.125in,height=\textheight]{../images//Consola2.png}
\end{center}
\end{frame}

\begin{frame}{Paneles de RStudio}
\phantomsection\label{paneles-de-rstudio}
\begin{columns}[T]
\begin{column}{0.5\textwidth}
\textbf{Panel superior izquierdo}

Editor de código --- aquí se escriben y editan los scripts

\textbf{Panel superior derecho}

\emph{Environment} e \emph{History} --- objetos activos e historial de
scripts
\end{column}

\begin{column}{0.5\textwidth}
\textbf{Panel inferior izquierdo}

Consola de R --- donde se ejecutan los códigos

\textbf{Panel inferior derecho}

\emph{Files} · \emph{Plots} · \emph{Packages} · \emph{Help} ·
\emph{Viewer}
\end{column}
\end{columns}
\end{frame}

\begin{frame}{\emph{Scripts}}
\phantomsection\label{scripts}
Una de las grandes ventajas de R es que los códigos e instrucciones se
guardan en \emph{scripts}, que se pueden editar y reutilizar.

\includegraphics[width=0.45\textwidth,height=\textheight]{../images//script.jpeg}

\includegraphics[width=0.45\textwidth,height=\textheight]{../images//script2.jpg}

Para iniciar un nuevo script: \emph{Archivo → Nuevo Archivo → R Script}
\end{frame}

\begin{frame}[fragile]{Estructura recomendada de un \emph{script}}
\phantomsection\label{estructura-recomendada-de-un-script}
\begin{Shaded}
\begin{Highlighting}[]
\CommentTok{\# Título / descripción del script}
\CommentTok{\# Autor: Stephanie · Fecha: 2026{-}05{-}19}

\FunctionTok{library}\NormalTok{(datasets)   }\CommentTok{\# primero: carga de paquetes}
\FunctionTok{data}\NormalTok{(iris)          }\CommentTok{\# segundo: carga de datos}

\FunctionTok{summary}\NormalTok{(iris)       }\CommentTok{\# análisis}
\FunctionTok{boxplot}\NormalTok{(iris)       }\CommentTok{\# gráficas}
\end{Highlighting}
\end{Shaded}

\begin{itemize}
\tightlist
\item
  El símbolo \texttt{\#} indica un comentario --- R lo ignora al
  ejecutar
\item
  Carguen siempre paquetes y datos \textbf{al principio} del script
\item
  Nombres de archivos: minúsculas, sin espacios → \texttt{mi\_script.R}
\end{itemize}
\end{frame}

\begin{frame}[fragile]{Abriendo un \emph{script}}
\phantomsection\label{abriendo-un-script}
Pueden abrir y ejecutar scripts en R con la función \texttt{source()}:

\begin{Shaded}
\begin{Highlighting}[]
\FunctionTok{source}\NormalTok{(}\StringTok{"C:/mi\_script.R"}\NormalTok{)}
\end{Highlighting}
\end{Shaded}

O bien: \emph{File → Open File} y buscar en sus carpetas.

\begin{block}{Comentarios}
\phantomsection\label{comentarios}
Siempre inicien con comentarios descriptivos. Usen \texttt{\#} para que
R los ignore:

\begin{Shaded}
\begin{Highlighting}[]
\CommentTok{\# Este es mi código}
\FunctionTok{data}\NormalTok{(iris)}
\CommentTok{\# Aquí termina el código}
\end{Highlighting}
\end{Shaded}
\end{block}
\end{frame}

\begin{frame}[fragile]{Cómo ejecutar comandos}
\phantomsection\label{cuxf3mo-ejecutar-comandos}
\begin{center}
\includegraphics[width=3.125in,height=\textheight]{../images//save.png}
\end{center}

\begin{itemize}
\tightlist
\item
  Ejecutar línea actual: \texttt{Ctrl\ +\ Enter} (Windows/Linux) ·
  \texttt{Cmd\ +\ Return} (Mac)
\item
  Ejecutar todo el script: botón \emph{Run} (esquina superior derecha
  del editor)
\item
  Guarden con un nombre descriptivo, ej. \texttt{mi\_script.R}
\end{itemize}
\end{frame}

\begin{frame}[fragile]{Ejemplo: corriendo el código}
\phantomsection\label{ejemplo-corriendo-el-cuxf3digo}
\begin{Shaded}
\begin{Highlighting}[]
\FunctionTok{library}\NormalTok{(datasets)}
\FunctionTok{data}\NormalTok{(iris)}
\FunctionTok{summary}\NormalTok{(iris)}
\FunctionTok{boxplot}\NormalTok{(iris)}
\end{Highlighting}
\end{Shaded}

\begin{center}
\includegraphics[width=3.33333in,height=1.66667in]{../images//code.png}
\end{center}
\end{frame}

\begin{frame}[fragile]{Directorio de trabajo}
\phantomsection\label{directorio-de-trabajo}
El \emph{directorio de trabajo} es el lugar en la computadora donde R
busca los archivos para importar y donde guarda los resultados.

\begin{Shaded}
\begin{Highlighting}[]
\FunctionTok{getwd}\NormalTok{()                          }\CommentTok{\# ver directorio actual}
\FunctionTok{setwd}\NormalTok{(}\StringTok{"/home/steph/Desktop/"}\NormalTok{)    }\CommentTok{\# establecer directorio}
\FunctionTok{list.files}\NormalTok{()                     }\CommentTok{\# archivos en el directorio}
\FunctionTok{list.dirs}\NormalTok{()                      }\CommentTok{\# carpetas en el directorio}
\end{Highlighting}
\end{Shaded}

\medskip

\begin{block}{Mejor práctica}
Usar \textbf{Proyectos de R} en lugar de \texttt{setwd()} — el directorio se configura automáticamente al abrir el proyecto.
\end{block}
\end{frame}

\begin{frame}[fragile]{Sesión de R}
\phantomsection\label{sesiuxf3n-de-r}
Los objetos y funciones de R se almacenan en la \textbf{memoria} de la
computadora. Cada vez que ejecutamos R, iniciamos una \emph{sesión}.

\begin{itemize}
\tightlist
\item
  Los objetos creados permanecen solo en esa sesión
\item
  Las sesiones no se comparten entre sí
\item
  Al cerrar R, puede guardarse todo en un archivo
  \textbf{\texttt{.Rdata}}
\end{itemize}

\begin{Shaded}
\begin{Highlighting}[]
\FunctionTok{ls}\NormalTok{()    }\CommentTok{\# lista todos los objetos en la sesión actual}
\end{Highlighting}
\end{Shaded}

\begin{alertblock}{No recomendado}
Guardar el espacio de trabajo completo genera confusión entre proyectos y ocupa mucho espacio en disco. Mejor exportar solo lo necesario.
\end{alertblock}
\end{frame}

\begin{frame}{\emph{Proyecto} de R (.Rproj)}
\phantomsection\label{proyecto-de-r-.rproj}
\begin{columns}[T]
\begin{column}{0.55\textwidth}
\begin{itemize}
\tightlist
\item
  Identifica todos los archivos y contenido asociados a un trabajo
\item
  Al crear un proyecto, las carpetas quedan \textbf{vinculadas
  directamente} a él
\item
  El directorio de trabajo se establece automáticamente
\item
  Un proyecto por cada curso, artículo o análisis
\end{itemize}
\end{column}

\begin{column}{0.45\textwidth}
\includegraphics[width=0.95\textwidth,height=\textheight]{../images//new_directory.png}

\includegraphics[width=0.95\textwidth,height=\textheight]{../images//new_directory2.png}
\end{column}
\end{columns}
\end{frame}

\begin{frame}{Crear un proyecto}
\phantomsection\label{crear-un-proyecto}
\begin{center}
\includegraphics[width=3.125in,height=\textheight]{../images//new_directory3.png}
\end{center}
\end{frame}

\begin{frame}[fragile]{Instalación de paquetes}
\phantomsection\label{instalaciuxf3n-de-paquetes}
\begin{itemize}
\tightlist
\item
  Muchos desarrolladores crean \emph{paquetes} con funcionalidades
  adicionales
\item
  Disponibles en \textbf{CRAN}, \textbf{GitHub} y \textbf{Bioconductor}
\end{itemize}

\begin{Shaded}
\begin{Highlighting}[]
\FunctionTok{install.packages}\NormalTok{(}\StringTok{"tidyverse"}\NormalTok{)   }\CommentTok{\# instalar desde CRAN}
\end{Highlighting}
\end{Shaded}

En RStudio: pestaña \emph{Packages → Install}. Para cargar en cada
sesión:

\begin{Shaded}
\begin{Highlighting}[]
\FunctionTok{library}\NormalTok{(tidyverse)}
\end{Highlighting}
\end{Shaded}

\begin{alertblock}{Importante}
Los paquetes se instalan \textbf{una sola vez}, pero se cargan \textbf{en cada sesión nueva}.
\end{alertblock}
\end{frame}

\begin{frame}[fragile]{Instalación: versiones de desarrollo y
dependencias}
\phantomsection\label{instalaciuxf3n-versiones-de-desarrollo-y-dependencias}
Para paquetes en GitHub (versión en desarrollo):

\begin{Shaded}
\begin{Highlighting}[]
\NormalTok{devtools}\SpecialCharTok{::}\FunctionTok{install\_github}\NormalTok{(}\StringTok{"rstudio/rmarkdown"}\NormalTok{)}
\end{Highlighting}
\end{Shaded}

Los \texttt{::} llaman la función de un paquete sin cargarlo
permanentemente en la sesión.

\medskip

\begin{block}{Dependencias}
Instalar \texttt{tidyverse} instala varios paquetes a la vez — tiene muchas \textbf{dependencias}. Al cargarlo con \texttt{library()}, también se cargan sus dependencias.
\end{block}
\end{frame}

\begin{frame}[fragile]{Tipos de objetos en R}
\phantomsection\label{tipos-de-objetos-en-r}
La información en R se estructura en forma de objetos, visibles en el
panel \emph{Environment}:

\begin{Shaded}
\begin{Highlighting}[]
\NormalTok{a     }\OtherTok{\textless{}{-}} \DecValTok{1}                                                 \CommentTok{\# escalar}
\NormalTok{letra }\OtherTok{\textless{}{-}} \StringTok{"a"}                                               \CommentTok{\# caracter}
\NormalTok{b     }\OtherTok{\textless{}{-}} \FunctionTok{c}\NormalTok{(}\DecValTok{1}\NormalTok{, }\DecValTok{2}\NormalTok{, }\DecValTok{3}\NormalTok{)                                        }\CommentTok{\# vector}
\NormalTok{c     }\OtherTok{\textless{}{-}} \FunctionTok{matrix}\NormalTok{(}\DecValTok{1}\SpecialCharTok{:}\DecValTok{10}\NormalTok{)                                      }\CommentTok{\# matriz}
\NormalTok{d     }\OtherTok{\textless{}{-}} \FunctionTok{data.frame}\NormalTok{(}\AttributeTok{Especie=}\FunctionTok{c}\NormalTok{(}\StringTok{"A"}\NormalTok{,}\StringTok{"B"}\NormalTok{), }\AttributeTok{Long=}\FunctionTok{c}\NormalTok{(}\DecValTok{1}\NormalTok{,}\DecValTok{2}\NormalTok{))       }\CommentTok{\# dataframe}
\NormalTok{e     }\OtherTok{\textless{}{-}} \FunctionTok{list}\NormalTok{(}\FunctionTok{c}\NormalTok{(}\DecValTok{1}\SpecialCharTok{:}\DecValTok{20}\NormalTok{), }\FunctionTok{c}\NormalTok{(}\DecValTok{1}\SpecialCharTok{:}\DecValTok{10}\NormalTok{))                            }\CommentTok{\# lista}
\end{Highlighting}
\end{Shaded}
\end{frame}

\begin{frame}[fragile]{Definiendo variables}
\phantomsection\label{definiendo-variables}
\begin{itemize}
\tightlist
\item
  Usamos \texttt{\textless{}-} para asignar valores a las variables
\item
  También puede usarse \texttt{=}, pero \textbf{no se recomienda}: en R
  el \texttt{=} implica igualdad lógica y puede causar confusiones
\end{itemize}

\begin{Shaded}
\begin{Highlighting}[]
\NormalTok{a }\OtherTok{\textless{}{-}} \DecValTok{1}
\NormalTok{a             }\CommentTok{\# evaluar imprime el valor}
\end{Highlighting}
\end{Shaded}

\begin{verbatim}
[1] 1
\end{verbatim}

\begin{Shaded}
\begin{Highlighting}[]
\FunctionTok{print}\NormalTok{(a)      }\CommentTok{\# forma explícita de imprimir}
\end{Highlighting}
\end{Shaded}

\begin{verbatim}
[1] 1
\end{verbatim}
\end{frame}

\begin{frame}[fragile]{Visualizando variables asignadas}
\phantomsection\label{visualizando-variables-asignadas}
Los objetos declarados aparecen en el panel \textbf{Environment}
(superior derecho). Desde ahí pueden verse el tipo y tamaño de cada
objeto, o hacer clic para explorar tablas y listas.

Si intentamos imprimir un objeto que no existe en la sesión, R devuelve
un error:

\texttt{Error:\ object\ \textquotesingle{}f\textquotesingle{}\ not\ found}

\begin{block}{Tip}
Con \texttt{ls()} obtenemos una lista con los nombres de todo lo guardado en la sesión actual.
\end{block}
\end{frame}

\begin{frame}[fragile]{Tips para nombrar variables}
\phantomsection\label{tips-para-nombrar-variables}
\begin{columns}[T]
\begin{column}{0.5\textwidth}
\textbf{Sí}

\begin{itemize}
\tightlist
\item
  Empezar con letra: \texttt{mi\_var}
\item
  Sin espacios: \texttt{longitud\_total}
\item
  Minúsculas + guión bajo \texttt{\_}
\item
  Nombres descriptivos y significativos
\end{itemize}
\end{column}

\begin{column}{0.5\textwidth}
\textbf{No}

\begin{itemize}
\tightlist
\item
  Números al inicio: \texttt{1especie}
\item
  Con espacios: \texttt{mi\ dato}
\item
  Guiones medios: \texttt{mi-dato}
\item
  Sobreescribir funciones: \texttt{mean\ \textless{}-\ 5}
\item
  Caracteres especiales: \texttt{-\ ;\ .\ @\ ?}
\end{itemize}
\end{column}
\end{columns}
\end{frame}

\begin{frame}[fragile]{Guardar el espacio de trabajo}
\phantomsection\label{guardar-el-espacio-de-trabajo}
Los objetos permanecen en el espacio de trabajo hasta que finaliza la
sesión. Al salir, R pregunta si desean guardarlo (archivo
\texttt{.Rdata}).

\begin{alertblock}{No recomendado}
Sin organización por proyectos, es muy difícil darle seguimiento a los objetos guardados y ocupa mucho espacio en disco.
\end{alertblock}

\begin{exampleblock}{Recomendado}
Un proyecto por cada trabajo o tarea, decidiendo explícitamente qué exportar y qué no.
\end{exampleblock}
\end{frame}

\begin{frame}[fragile]{Exportar objetos de R}
\phantomsection\label{exportar-objetos-de-r}
\begin{Shaded}
\begin{Highlighting}[]
\FunctionTok{library}\NormalTok{(readr)}
\FunctionTok{write\_tsv}\NormalTok{(d, }\StringTok{"data.tsv"}\NormalTok{)          }\CommentTok{\# separado por tabs}
\FunctionTok{write.table}\NormalTok{(d, }\StringTok{"data.txt"}\NormalTok{, }\AttributeTok{sep=}\StringTok{"}\SpecialCharTok{\textbackslash{}t}\StringTok{"}\NormalTok{)}
\FunctionTok{write\_csv}\NormalTok{(d, }\StringTok{"data.csv"}\NormalTok{)          }\CommentTok{\# separado por comas}
\end{Highlighting}
\end{Shaded}
\end{frame}

\begin{frame}[fragile]{Exportar objetos: formato nativo de R}
\phantomsection\label{exportar-objetos-formato-nativo-de-r}
Para guardar objetos conservando su estructura en R:

\begin{Shaded}
\begin{Highlighting}[]
\FunctionTok{saveRDS}\NormalTok{(d, }\StringTok{"data.RDS"}\NormalTok{)}
\end{Highlighting}
\end{Shaded}

Para cargar el objeto nuevamente:

\begin{Shaded}
\begin{Highlighting}[]
\FunctionTok{readRDS}\NormalTok{(}\StringTok{"data.RDS"}\NormalTok{)}
\end{Highlighting}
\end{Shaded}

\begin{block}{Ventaja de RDS}
A diferencia de los formatos de texto, \texttt{saveRDS} conserva la clase y estructura del objeto (factores, listas, etc.)
\end{block}
\end{frame}

\begin{frame}[fragile]{Tipos de datos}
\phantomsection\label{tipos-de-datos}
La función \texttt{class()} determina qué tipo de objeto tenemos:

\begin{Shaded}
\begin{Highlighting}[]
\NormalTok{a }\OtherTok{\textless{}{-}} \DecValTok{5}
\FunctionTok{class}\NormalTok{(a)}
\end{Highlighting}
\end{Shaded}

\begin{verbatim}
[1] "numeric"
\end{verbatim}

\begin{longtable}[]{@{}lll@{}}
\toprule\noalign{}
Tipo & En R & Ejemplo \\
\midrule\noalign{}
\endhead
Numérico & \texttt{numeric} & \texttt{5.1} \\
Entero & \texttt{integer} & \texttt{4L} \\
Texto & \texttt{character} & \texttt{"hola"} \\
Factor & \texttt{factor} & Bajo / Medio / Alto \\
Lógico & \texttt{logical} & \texttt{TRUE}, \texttt{FALSE} \\
Perdido & \texttt{NA} & \texttt{NA} \\
Vacío & \texttt{NULL} & \texttt{NULL} \\
\bottomrule\noalign{}
\end{longtable}
\end{frame}

\begin{frame}{Para tener en cuenta\ldots{}}
\phantomsection\label{para-tener-en-cuenta}
\begin{block}{character}
Siempre rodeado de comillas simples o dobles. Es el tipo más flexible — puede contener letras, números, espacios y símbolos especiales.
\end{block}

\begin{alertblock}{factor}
Dato numérico representado con etiquetas (niveles). Los niveles se ordenan \textbf{alfabéticamente} por defecto. Puede causar confusión: a veces se comporta como carácter y otras veces no — es fuente común de errores.
\end{alertblock}

\begin{exampleblock}{NA vs NULL}
\texttt{NULL} aparece cuando R intenta recuperar un dato y no encuentra nada. \texttt{NA} representa explícitamente un \textbf{dato faltante o perdido} y puede surgir de operaciones que no pudieron completarse.
\end{exampleblock}
\end{frame}

\begin{frame}[fragile]{Coerción de datos}
\phantomsection\label{coerciuxf3n-de-datos}
En R los datos pueden ser \emph{coercionados} --- forzados a
transformarse de un tipo a otro:

\begin{columns}[T]
\begin{column}{0.5\textwidth}
\begin{Shaded}
\begin{Highlighting}[]
\NormalTok{a }\OtherTok{\textless{}{-}} \DecValTok{5}
\FunctionTok{as.character}\NormalTok{(a)}
\end{Highlighting}
\end{Shaded}

\begin{verbatim}
[1] "5"
\end{verbatim}

\begin{Shaded}
\begin{Highlighting}[]
\FunctionTok{as.factor}\NormalTok{(}\StringTok{"medio"}\NormalTok{)}
\end{Highlighting}
\end{Shaded}

\begin{verbatim}
[1] medio
Levels: medio
\end{verbatim}
\end{column}

\begin{column}{0.5\textwidth}
\begin{longtable}[]{@{}ll@{}}
\toprule\noalign{}
Función & Convierte a \\
\midrule\noalign{}
\endhead
\texttt{as.integer()} & Entero \\
\texttt{as.numeric()} & Numérico \\
\texttt{as.character()} & Texto \\
\texttt{as.factor()} & Factor \\
\texttt{as.logical()} & Lógico \\
\bottomrule\noalign{}
\end{longtable}
\end{column}
\end{columns}
\end{frame}

\begin{frame}[fragile]{Tipos de estructura de datos}
\phantomsection\label{tipos-de-estructura-de-datos}
Los datos se estructuran de distintas formas. \texttt{class()} también
identifica la estructura:

\begin{columns}[T]
\begin{column}{0.5\textwidth}
\textbf{Homogéneas} (un solo tipo)

\begin{itemize}
\tightlist
\item
  \textbf{Vector} --- 1 dimensión
\item
  \textbf{Matriz} --- 2 dimensiones
\end{itemize}
\end{column}

\begin{column}{0.5\textwidth}
\textbf{Heterogéneas} (tipos mixtos)

\begin{itemize}
\tightlist
\item
  \textbf{Lista} --- 1 dimensión
\item
  \textbf{Data frame} --- 2 dimensiones
\end{itemize}
\end{column}
\end{columns}
\end{frame}

\begin{frame}[fragile]{Vectores}
\phantomsection\label{vectores}
Colecciones de uno o más datos del \textbf{mismo tipo}. No es posible
mezclar tipos:

\begin{Shaded}
\begin{Highlighting}[]
\NormalTok{colores }\OtherTok{\textless{}{-}} \FunctionTok{c}\NormalTok{(}\StringTok{"red"}\NormalTok{, }\StringTok{"black"}\NormalTok{, }\StringTok{"blue"}\NormalTok{)}
\FunctionTok{is.vector}\NormalTok{(colores)}
\end{Highlighting}
\end{Shaded}

\begin{verbatim}
[1] TRUE
\end{verbatim}

\begin{Shaded}
\begin{Highlighting}[]
\FunctionTok{class}\NormalTok{(colores)}
\end{Highlighting}
\end{Shaded}

\begin{verbatim}
[1] "character"
\end{verbatim}
\end{frame}

\begin{frame}[fragile]{Vectorización}
\phantomsection\label{vectorizaciuxf3n}
Las operaciones aritméticas y relacionales se aplican a \textbf{cada
elemento} del vector automáticamente:

\begin{Shaded}
\begin{Highlighting}[]
\NormalTok{un\_vector }\OtherTok{\textless{}{-}} \FunctionTok{c}\NormalTok{(}\DecValTok{1}\SpecialCharTok{:}\DecValTok{10}\NormalTok{)}
\NormalTok{un\_vector }\SpecialCharTok{*} \DecValTok{10}
\end{Highlighting}
\end{Shaded}

\begin{verbatim}
 [1]  10  20  30  40  50  60  70  80  90 100
\end{verbatim}

\begin{Shaded}
\begin{Highlighting}[]
\NormalTok{un\_vector }\SpecialCharTok{+} \DecValTok{1}
\end{Highlighting}
\end{Shaded}

\begin{verbatim}
 [1]  2  3  4  5  6  7  8  9 10 11
\end{verbatim}

\begin{Shaded}
\begin{Highlighting}[]
\NormalTok{un\_vector }\SpecialCharTok{+}\NormalTok{ un\_vector}
\end{Highlighting}
\end{Shaded}

\begin{verbatim}
 [1]  2  4  6  8 10 12 14 16 18 20
\end{verbatim}
\end{frame}

\begin{frame}[fragile]{Matrices y arreglos}
\phantomsection\label{matrices-y-arreglos}
\begin{itemize}
\tightlist
\item
  Vectores multidimensionales --- deben contener un \textbf{solo tipo}
  de datos
\item
  Una matriz tiene \textbf{2 dimensiones}; un arreglo (\emph{array})
  puede tener \texttt{n} dimensiones
\end{itemize}

\begin{Shaded}
\begin{Highlighting}[]
\NormalTok{vect  }\OtherTok{\textless{}{-}} \DecValTok{1}\SpecialCharTok{:}\DecValTok{20}
\NormalTok{matr  }\OtherTok{\textless{}{-}} \FunctionTok{matrix}\NormalTok{(}\DecValTok{1}\SpecialCharTok{:}\DecValTok{20}\NormalTok{)}
\NormalTok{matri }\OtherTok{\textless{}{-}} \FunctionTok{matrix}\NormalTok{(}\DecValTok{1}\SpecialCharTok{:}\DecValTok{20}\NormalTok{, }\AttributeTok{nrow=}\DecValTok{5}\NormalTok{, }\AttributeTok{ncol=}\DecValTok{4}\NormalTok{)}
\FunctionTok{dim}\NormalTok{(matri)}
\end{Highlighting}
\end{Shaded}

\begin{verbatim}
[1] 5 4
\end{verbatim}

\begin{Shaded}
\begin{Highlighting}[]
\NormalTok{matri}
\end{Highlighting}
\end{Shaded}

\begin{verbatim}
     [,1] [,2] [,3] [,4]
[1,]    1    6   11   16
[2,]    2    7   12   17
[3,]    3    8   13   18
[4,]    4    9   14   19
[5,]    5   10   15   20
\end{verbatim}
\end{frame}

\begin{frame}[fragile]{Operaciones con matrices}
\phantomsection\label{operaciones-con-matrices}
Con \texttt{dim()} conocemos las dimensiones (filas × columnas):

\begin{Shaded}
\begin{Highlighting}[]
\NormalTok{matri }\SpecialCharTok{*} \DecValTok{2}
\end{Highlighting}
\end{Shaded}

\begin{verbatim}
     [,1] [,2] [,3] [,4]
[1,]    2   12   22   32
[2,]    4   14   24   34
[3,]    6   16   26   36
[4,]    8   18   28   38
[5,]   10   20   30   40
\end{verbatim}
\end{frame}

\begin{frame}[fragile]{Transponer una matriz}
\phantomsection\label{transponer-una-matriz}
Podemos rotar la matriz 90° con la función \texttt{t()}:

\begin{Shaded}
\begin{Highlighting}[]
\FunctionTok{t}\NormalTok{(matri)}
\end{Highlighting}
\end{Shaded}

\begin{verbatim}
     [,1] [,2] [,3] [,4] [,5]
[1,]    1    2    3    4    5
[2,]    6    7    8    9   10
[3,]   11   12   13   14   15
[4,]   16   17   18   19   20
\end{verbatim}
\end{frame}

\begin{frame}[fragile]{Listas}
\phantomsection\label{listas}
Las listas son \textbf{heterogéneas}: cada elemento puede ser de
distinto tipo o clase:

\begin{Shaded}
\begin{Highlighting}[]
\NormalTok{un\_vector }\OtherTok{\textless{}{-}} \DecValTok{1}\SpecialCharTok{:}\DecValTok{20}
\NormalTok{una\_matriz }\OtherTok{\textless{}{-}} \FunctionTok{matrix}\NormalTok{(}\DecValTok{1}\SpecialCharTok{:}\DecValTok{20}\NormalTok{, }\AttributeTok{nrow=}\DecValTok{2}\NormalTok{)}
\NormalTok{una\_df }\OtherTok{\textless{}{-}} \FunctionTok{data.frame}\NormalTok{(}\StringTok{"numeros"}\OtherTok{=}\DecValTok{1}\SpecialCharTok{:}\DecValTok{3}\NormalTok{, }\StringTok{"letras"}\OtherTok{=}\FunctionTok{c}\NormalTok{(}\StringTok{"a"}\NormalTok{,}\StringTok{"b"}\NormalTok{,}\StringTok{"c"}\NormalTok{))}
\NormalTok{una\_lista }\OtherTok{\textless{}{-}} \FunctionTok{list}\NormalTok{(}\StringTok{"vec"}\OtherTok{=}\NormalTok{un\_vector, }\StringTok{"mat"}\OtherTok{=}\NormalTok{una\_matriz, }\StringTok{"df"}\OtherTok{=}\NormalTok{una\_df)}
\FunctionTok{str}\NormalTok{(una\_lista)}
\end{Highlighting}
\end{Shaded}

\begin{verbatim}
List of 3
 $ vec: int [1:20] 1 2 3 4 5 6 7 8 9 10 ...
 $ mat: int [1:2, 1:10] 1 2 3 4 5 6 7 8 9 10 ...
 $ df :'data.frame':    3 obs. of  2 variables:
  ..$ numeros: int [1:3] 1 2 3
  ..$ letras : chr [1:3] "a" "b" "c"
\end{verbatim}
\end{frame}

\begin{frame}[fragile]{Coercionando entre estructuras}
\phantomsection\label{coercionando-entre-estructuras}
\begin{Shaded}
\begin{Highlighting}[]
\NormalTok{df\_transpuesta }\OtherTok{\textless{}{-}} \FunctionTok{t}\NormalTok{(una\_df)}
\FunctionTok{class}\NormalTok{(df\_transpuesta)}
\end{Highlighting}
\end{Shaded}

\begin{verbatim}
[1] "matrix" "array" 
\end{verbatim}

La función \texttt{t()} cambia la clase a matrix. Para conservar el
dataframe:

\begin{Shaded}
\begin{Highlighting}[]
\NormalTok{df\_transpuesta }\OtherTok{\textless{}{-}} \FunctionTok{as.data.frame}\NormalTok{(}\FunctionTok{t}\NormalTok{(una\_df))}
\FunctionTok{class}\NormalTok{(df\_transpuesta)}
\end{Highlighting}
\end{Shaded}

\begin{verbatim}
[1] "data.frame"
\end{verbatim}
\end{frame}

\begin{frame}[fragile]{Funciones}
\phantomsection\label{funciones}
\begin{itemize}
\tightlist
\item
  A lo largo del curso ya usamos varias funciones: \texttt{library()},
  \texttt{write\_csv()}, \texttt{ls()}\ldots{}
\item
  La \textbf{sintaxis} de R requiere paréntesis para evaluar funciones
\item
  R incluye muchas funciones por defecto; otras vienen en paquetes
\end{itemize}

\begin{Shaded}
\begin{Highlighting}[]
\FunctionTok{log}\NormalTok{(a)          }\CommentTok{\# logaritmo natural de a}
\end{Highlighting}
\end{Shaded}

\begin{verbatim}
[1] 1.609438
\end{verbatim}

\begin{Shaded}
\begin{Highlighting}[]
\FunctionTok{sqrt}\NormalTok{(a)         }\CommentTok{\# raíz cuadrada de a}
\end{Highlighting}
\end{Shaded}

\begin{verbatim}
[1] 2.236068
\end{verbatim}

Para ver los argumentos sin abrir la ayuda:

\begin{Shaded}
\begin{Highlighting}[]
\FunctionTok{args}\NormalTok{(log)}
\end{Highlighting}
\end{Shaded}

\begin{verbatim}
function (x, base = exp(1)) 
NULL
\end{verbatim}
\end{frame}

\begin{frame}[fragile]{Funciones: argumentos}
\phantomsection\label{funciones-argumentos}
\begin{Shaded}
\begin{Highlighting}[]
\FunctionTok{help}\NormalTok{(}\StringTok{"log"}\NormalTok{)}
\NormalTok{?log}
\end{Highlighting}
\end{Shaded}

\begin{itemize}
\tightlist
\item
  Argumentos con \texttt{=} en la ayuda son \textbf{opcionales} ---
  tienen valor por defecto
\item
  Ejemplo: \texttt{log} usa \texttt{base\ =\ exp(1)} por defecto
  (logaritmo natural)
\end{itemize}

\begin{Shaded}
\begin{Highlighting}[]
\FunctionTok{log}\NormalTok{(}\AttributeTok{x=}\DecValTok{8}\NormalTok{, }\AttributeTok{base=}\DecValTok{2}\NormalTok{)   }\CommentTok{\# argumentos por nombre}
\end{Highlighting}
\end{Shaded}

\begin{verbatim}
[1] 3
\end{verbatim}

\begin{Shaded}
\begin{Highlighting}[]
\FunctionTok{log}\NormalTok{(}\DecValTok{8}\NormalTok{, }\DecValTok{2}\NormalTok{)          }\CommentTok{\# argumentos por posición — mismo resultado}
\end{Highlighting}
\end{Shaded}

\begin{verbatim}
[1] 3
\end{verbatim}

En las funciones sí se usa \texttt{=}, no \texttt{\textless{}-}, para
especificar argumentos.
\end{frame}

\begin{frame}[fragile]{Funciones: operadores aritméticos}
\phantomsection\label{funciones-operadores-aritmuxe9ticos}
Los operadores son funciones pero sin paréntesis:

\begin{Shaded}
\begin{Highlighting}[]
\NormalTok{a }\SpecialCharTok{+}\NormalTok{ a}
\end{Highlighting}
\end{Shaded}

\begin{verbatim}
[1] 10
\end{verbatim}

\begin{Shaded}
\begin{Highlighting}[]
\NormalTok{a }\SpecialCharTok{{-}} \DecValTok{2}
\end{Highlighting}
\end{Shaded}

\begin{verbatim}
[1] 3
\end{verbatim}

\begin{Shaded}
\begin{Highlighting}[]
\DecValTok{1} \SpecialCharTok{*}\NormalTok{ pi}
\end{Highlighting}
\end{Shaded}

\begin{verbatim}
[1] 3.141593
\end{verbatim}

\begin{Shaded}
\begin{Highlighting}[]
\DecValTok{2} \SpecialCharTok{/} \DecValTok{3}
\end{Highlighting}
\end{Shaded}

\begin{verbatim}
[1] 0.6666667
\end{verbatim}

\begin{Shaded}
\begin{Highlighting}[]
\DecValTok{4} \SpecialCharTok{\^{}}\NormalTok{ a}
\end{Highlighting}
\end{Shaded}

\begin{verbatim}
[1] 1024
\end{verbatim}
\end{frame}

\begin{frame}[fragile]{Funciones: crear las tuyas (1)}
\phantomsection\label{funciones-crear-las-tuyas-1}
Es posible declarar funciones propias. Por ejemplo, si quisiéramos
definir el promedio:

\begin{Shaded}
\begin{Highlighting}[]
\NormalTok{average }\OtherTok{\textless{}{-}} \ControlFlowTok{function}\NormalTok{(x) \{ }\FunctionTok{sum}\NormalTok{(x) }\SpecialCharTok{/} \FunctionTok{length}\NormalTok{(x) \}}
\NormalTok{x }\OtherTok{\textless{}{-}} \DecValTok{1}\SpecialCharTok{:}\DecValTok{100}
\FunctionTok{average}\NormalTok{(x)}
\end{Highlighting}
\end{Shaded}

\begin{verbatim}
[1] 50.5
\end{verbatim}

\begin{Shaded}
\begin{Highlighting}[]
\FunctionTok{mean}\NormalTok{(x)       }\CommentTok{\# comparamos con la función de R base}
\end{Highlighting}
\end{Shaded}

\begin{verbatim}
[1] 50.5
\end{verbatim}

Al compararlas obtenemos el mismo resultado.
\end{frame}

\begin{frame}[fragile]{Funciones: crear las tuyas (2)}
\phantomsection\label{funciones-crear-las-tuyas-2}
La estructura general de una función en R es:

\begin{Shaded}
\begin{Highlighting}[]
\NormalTok{nombre }\OtherTok{\textless{}{-}} \ControlFlowTok{function}\NormalTok{(argumentos) \{}
\NormalTok{  operaciones}
\NormalTok{\}}
\end{Highlighting}
\end{Shaded}

\begin{itemize}
\tightlist
\item
  Primero se declara el nombre y los argumentos
\item
  Luego las operaciones que realizará
\item
  Lo importante: hemos declarado un \textbf{nuevo objeto} de tipo
  función en nuestra sesión
\end{itemize}
\end{frame}




\end{document}
